\subsection{Hipercohomología de los complejos de módulos sobre un preesquema}
\label{subsection:III.6.2}

\begin{env}[6.2.1]
Sean $X$ un preesquema y
$\mathscr K^\bullet=(\mathscr K^i)_{i\in\mathbf Z}$ un complejo de
$\mathscr O_X$-módulos cuya diferencial es de grado $+1$.
Recordemos que, para todo morfismo $f:X\to Y$ de preesquemas, se han definido
(0, 12.4.1) los $\mathscr O_Y$-módulos de hipercohomología
$\mathscr H^n(f,\mathscr K^\bullet)$ (denotados también por
$\mathscr H_f^n(\mathscr K^\bullet)$ o
$\mathrm R^n f_*(\mathscr K^\bullet)$) para todo $n\in\mathbf Z$;
la hipercohomología $\mathscr H^\bullet(f,\mathscr K^\bullet)$ es el objeto de
convergencia de los dos funtores espectrales
$'\mathscr E(f,\mathscr K^\bullet)$ y
$''\mathscr E(f,\mathscr K^\bullet)$, cuyos términos $\mathscr E_2$ vienen
dados por
\begin{align}
  '\mathscr E_2^{pq}
    &=\mathscr H^p\bigl(\mathscr H^q(f,\mathscr K^\bullet)\bigr),
    \tag{6.2.1.1}\label{III.6.2.1.1}\\
  ''\mathscr E_2^{pq}
    &=\mathscr H^p\bigl(f,\mathscr H^q(\mathscr K^\bullet)\bigr)
      =\mathrm R^p f_*\bigl(\mathscr H^q(\mathscr K^\bullet)\bigr).
    \tag{6.2.1.2}\label{III.6.2.1.2}
\end{align}

\oldpage[III]{139}

donde $\mathscr H^q(f,\mathscr K^\bullet)$ es el complejo cuyo componente de
grado $i$ es
$\mathscr H^q(f,\mathscr K^i)=\mathrm R^q f_*(\mathscr K^i)$ (loc. cit.).
Recordemos también que, cuando $Y$ se reduce a un punto, se denota la
hipercohomología correspondiente por $\mathbf H^\bullet(X,\mathscr K^\bullet)$
(la cual está formada por módulos sobre $\Gamma(X,\mathscr O_X)$ independientes
del preesquema puntual $Y$ considerado); cuando $Y=X$ y $f=1_X$, se tiene
$\mathscr H^n(f,\mathscr K^\bullet)=\mathscr H^n(\mathscr K^\bullet)$
(cohomología del complejo $\mathscr K^\bullet$); cuando
$\mathscr K^i=0$, salvo para $i=i_0$, se tiene
\[
  \mathscr H^n(f,\mathscr K^\bullet)
    =\mathrm R^{n-i_0}f_*(\mathscr K^{i_0}).
\]
La sucesión espectral $'\mathscr E(f,\mathscr K^\bullet)$ es siempre
regular; las dos sucesiones espectrales son birregulares cuando
$\mathscr K^\bullet$ está acotado inferiormente (0, 12.4.1).

Toda equivalencia homotópica $h:\mathscr K^\bullet\to\mathscr K'^\bullet$ de complejos de
$\mathscr O_X$-módulos (6.1.4) da lugar a un isomorfismo
$\mathscr H^\bullet(f,\mathscr K^\bullet)\xrightarrow{\sim}
\mathscr H^\bullet(f,\mathscr K'^\bullet)$ en hipercohomología. Lo mismo ocurre
si se supone solamente que
$\mathscr H^\bullet(h):\mathscr H^\bullet(\mathscr K^\bullet)\to
\mathscr H^\bullet(\mathscr K'^\bullet)$ es un isomorfismo y que
$\mathscr K^\bullet$ y $\mathscr K'^\bullet$ están acotados inferiormente,
como se deduce inmediatamente de (0, 11.1.5), aplicado a la sucesión espectral
(6.2.1.2) y a su análoga para $\mathscr K'^\bullet$. Por último, para todo
cubrimiento abierto $\mathfrak U=(U_\alpha)$ de $X$, se ha definido también
(0, 12.4.5) la hipercohomología
$\mathbf H^\bullet(\mathfrak U,\mathscr K^\bullet)$ como la cohomología del
bicomplejo $C^\bullet(\mathfrak U,\mathscr K^\bullet)$ (cuyo componente de
índices $(i,j)$ es, por definición,
$C^i(\mathfrak U,\mathscr K^j)$); los
$\mathbf H^n(\mathfrak U,\mathscr K^\bullet)$ son también módulos sobre
$\Gamma(X,\mathscr O_X)$.
\end{env}

\begin{proposition}[6.2.2]
Sean $X$ un esquema y $\mathfrak U=(U_\alpha)$ un cubrimiento de $X$ por
abiertos afines. Para todo complejo $\mathscr K^\bullet$ de
$\mathscr O_X$-módulos cuasicoherentes, los módulos de hipercohomología
$\mathbf H^\bullet(X,\mathscr K^\bullet)$ y
$\mathbf H^\bullet(\mathfrak U,\mathscr K^\bullet)$ son canónicamente
isomorfos.
\end{proposition}

\begin{proof}
En efecto, toda intersección finita $V$ de abiertos del cubrimiento
$\mathfrak U$ es afín (I, 5.5.6), por lo que
$\mathrm H^q(V,\mathscr K^i)=0$ para todo $i$ y todo $q>0$ (1.3.1); la
proposición es, por tanto, un caso particular de (0, 12.4.7).
\end{proof}

\begin{proposition}[6.2.3]
Sea $f:X\to Y$ un morfismo cuasicompacto y separado de preesquemas. Para todo
complejo $\mathscr K^\bullet$ de $\mathscr O_X$-módulos cuasicoherentes, los
$\mathscr O_Y$-módulos $\mathscr H^n(f,\mathscr K^\bullet)$ son
cuasicoherentes.
\end{proposition}

\begin{proof}
Como los $\mathscr H^q(f,\mathscr K^i)=\mathrm R^q f_*(\mathscr K^i)$ son
$\mathscr O_Y$-módulos cuasicoherentes (1.4.10), lo mismo ocurre con
$'\mathscr E_2^{pq}$, que, según (6.2.1.1), es cociente del núcleo de un
homomorfismo de módulos cuasicoherentes por la imagen de un homomorfismo de
esta clase (I, 4.1.1). Por la misma razón, todos los
$\mathscr O_Y$-módulos $'\mathscr E_r^{pq}$,
$B_k('\mathscr E_2^{pq})$, $Z_k('\mathscr E_2^{pq})$ de la primera sucesión
espectral son cuasicoherentes. La regularidad de la sucesión espectral
$'\mathscr E(f,\mathscr K^\bullet)$ implica que
$Z_\infty('\mathscr E_2^{pq})$ es igual a uno de los
$Z_k('\mathscr E_2^{pq})$, luego es cuasicoherente, y lo mismo ocurre con
$B_\infty('\mathscr E_2^{pq})=\varinjlim_k B_k('\mathscr E_2^{pq})$
(0, 11.2.4 e I, 4.1.1); por tanto, los $'\mathscr E_\infty^{pq}$ son también
cuasicoherentes. Como la sucesión espectral precedente es regular, la
filtración de los $F^p(\mathscr H^n(f,\mathscr K^\bullet))$ es discreta y
exhaustiva; dicho de otro modo, el $\mathscr O_Y$-módulo
$\mathscr H^n(f,\mathscr K^\bullet)$ es la unión de una sucesión creciente
$(\mathscr G_k)_{k\geq0}$ de $\mathscr O_Y$-módulos tal que
$\mathscr G_0=0$ y cada $\mathscr G_k/\mathscr G_{k-1}$ es igual a uno de los
$\mathscr O_Y$-módulos $'\mathscr E_\infty^{pq}$ y, por tanto, es
cuasicoherente. Por inducción sobre $k$, se deduce que los $\mathscr G_k$ son
cuasicoherentes (1.4.17), y, como
$\mathscr H^n(f,\mathscr K^\bullet)=\varinjlim_k\mathscr G_k$, queda demostrada
la proposición (I, 4.1.1).
\end{proof}

\oldpage[III]{140}

\begin{corollary}[6.2.4]
Bajo las hipótesis de (6.2.3), para todo abierto afín $V$ de $Y$, el
homomorfismo canónico
\[
  \mathbf H^n(f^{-1}(V),\mathscr K^\bullet)
  \longrightarrow
  \Gamma(V,\mathscr H^n(f,\mathscr K^\bullet))
  \tag{6.2.4.1}\label{III.6.2.4.1}
\]
es biyectivo para todo $n\in\mathbf Z$.
\end{corollary}

\begin{proof}
La demostración es la misma que la de (1.4.11), utilizando (6.2.2),
reemplazando $\mathscr F$ por $\mathscr K^\bullet$,
$\mathscr K^\bullet$ por
$f_*(C^\bullet(\mathfrak U,\mathscr K^\bullet))$,
$\mathscr H^\bullet(\mathscr K^\bullet)$ por
$\mathscr H^\bullet(f,\mathscr K^\bullet)$, y observando que este último es un
$\mathscr O_Y$-módulo cuasicoherente por (6.2.3).
\end{proof}

\begin{proposition}[6.2.5]
Sean $Y$ un preesquema localmente noetheriano, $f:X\to Y$ un morfismo propio
y $\mathscr K^\bullet$ un complejo de $\mathscr O_X$-módulos tal que los
$\mathscr O_X$-módulos $\mathscr H^q(\mathscr K^\bullet)$ sean coherentes.
Entonces los $\mathscr O_Y$-módulos $\mathscr H^n(f,\mathscr K^\bullet)$ son
coherentes.
\end{proposition}

\begin{proof}
Como la cuestión es local sobre $Y$, podemos limitarnos al caso en que $Y$ es
noetheriano y afín; se trata entonces, en virtud de (6.2.4), de demostrar que
los $\mathbf H^n(X,\mathscr K^\bullet)$ son
$\Gamma(Y,\mathscr O_Y)$-módulos de tipo finito. Se tiene entonces
$\mathbf H^\bullet(X,\mathscr K^\bullet)=
\mathbf H^\bullet(\mathfrak U,\mathscr K^\bullet)$ (6.2.2), donde se puede
suponer que $\mathfrak U$ es finito, puesto que $X$ es cuasicompacto. Como las
cocadenas de cada complejo $C^\bullet(\mathfrak U,\mathscr K^j)$ son
alternadas por definición, existe un entero $r>0$ tal que
$C^i(\mathfrak U,\mathscr K^j)=0$ para $i<0$ y $i>r$; se concluye de ello
(0, 11.3.3) que las dos sucesiones espectrales del bicomplejo
$C^\bullet(\mathfrak U,\mathscr K^\bullet)$ son birregulares. Como las
intersecciones de los miembros de $\mathfrak U$ son abiertos afines
(I, 5.5.6), cada funtor
$\mathscr F\mapsto C^i(\mathfrak U,\mathscr F)$ es exacto en la categoría de
los $\mathscr O_X$-módulos cuasicoherentes; por consiguiente,
$\mathrm H_I^q(C^i(\mathfrak U,\mathscr K^\bullet))
=C^i(\mathfrak U,\mathscr H^q(\mathscr K^\bullet))$, y los términos
$E_2$ de la segunda sucesión espectral de
$C^\bullet(\mathfrak U,\mathscr K^\bullet)$ vienen dados (0, 11.3.2) por
\[
  ''E_2^{pq}
  =\mathrm H^p(C^\bullet(\mathfrak U,\mathscr H^q(\mathscr K^\bullet)))
  =\mathbf H^p(\mathfrak U,\mathscr H^q(\mathscr K^\bullet))
  =\mathrm H^p(X,\mathscr H^q(\mathscr K^\bullet))
\]
en virtud de (1.4.1); dado que $f$ es propio, estos son
$\Gamma(Y,\mathscr O_Y)$-módulos de tipo finito (3.2.1). Puesto que la sucesión
espectral $''E(C^\bullet(\mathfrak U,\mathscr K^\bullet))$ es birregular, se
deduce que los $\mathbf H^n(X,\mathscr K^\bullet)$ son
$\Gamma(Y,\mathscr O_Y)$-módulos de tipo finito (0, 11.1.8).

A fortiori, si $\mathscr K^\bullet$ es un complejo de
$\mathscr O_X$-módulos coherentes, los $\mathscr O_Y$-módulos
$\mathscr H^n(f,\mathscr K^\bullet)$ son coherentes bajo las hipótesis de
(6.2.5) relativas a $Y$ y $f$ (0\textsubscript{I}, 5.3.4).
\end{proof}

\begin{env}[6.2.6]
La hipercohomología $\mathscr H^\bullet(f,\mathscr K^\bullet)$ es un funtor
cohomológico en la categoría de los complejos de $\mathscr O_X$-módulos
acotados inferiormente (0, 12.4.4). Es un funtor cohomológico en la categoría
de todos los complejos de $\mathscr O_X$-módulos cuando el morfismo $f$ es
cuasicompacto y el espacio subyacente a $X$ es localmente noetheriano: en
efecto, de (G, II, 3.10.1) se deduce entonces que $f_*$ conmuta con los límites
inductivos (pues la cuestión es local sobre $Y$), y puede aplicarse
(0, 11.5.2).

Por último, si $f$ es separado, $\mathscr H^\bullet(f,\mathscr K^\bullet)$ es
un funtor cohomológico en la categoría de los complejos de
$\mathscr O_X$-módulos cuasicoherentes. Esto es inmediato cuando $Y$ es
afín, pues entonces $X$ es un esquema; por tanto, en virtud del isomorfismo
canónico (6.2.2), todo se reduce a ver que
$\mathscr K^\bullet\mapsto\mathbf H^\bullet(\mathfrak U,\mathscr K^\bullet)$
es un funtor cohomológico en la categoría de los complejos de
$\mathscr O_X$-módulos cuasicoherentes, lo cual es inmediato, ya que el
funtor $\mathscr K^\bullet\mapsto
C^\bullet(\mathfrak U,\mathscr K^\bullet)$ es exacto en esta categoría
(I, 1.3.7). En el caso general, para todo abierto afín $V$

\oldpage[III]{141}

de $Y$, $f^{-1}(V)$ es un esquema y, para aplicar lo que precede, basta
verificar que, para una sucesión exacta
\[
  0\longrightarrow\mathscr K'^\bullet
  \longrightarrow\mathscr K^\bullet
  \longrightarrow\mathscr K''^\bullet
  \longrightarrow0
\]
de complejos de $\mathscr O_X$-módulos cuasicoherentes, el homomorfismo
\[
  \partial:
  \mathscr H^n(f,\mathscr K''^\bullet)|V
  \longrightarrow
  \mathscr H^{n+1}(f,\mathscr K'^\bullet)|V
\]
no depende del cubrimiento abierto afín $\mathfrak U$ de $f^{-1}(V)$ utilizado
para definirlo. Pero esto se deduce de que, si $\mathfrak U'$ es un
cubrimiento abierto afín más fino que $\mathfrak U$, el diagrama
\[
  \xymatrix@C=34pt@R=11pt{
    \mathbf H^\bullet(\mathfrak U,\mathscr K^\bullet)
      \ar[rrd]^-{\sim}\ar[dd]^{\iota}
    && \\
    &&
    \mathbf H^\bullet(f^{-1}(V),\mathscr K^\bullet)
    \\
    \mathbf H^\bullet(\mathfrak U',\mathscr K^\bullet)
      \ar[rru]_-{\sim}
    &&
  }
\]
de isomorfismos canónicos es conmutativo, al igual que el diagrama
\[
  \xymatrix@C=46pt@R=25pt{
    \mathbf H^n(\mathfrak U,\mathscr K''^\bullet)
      \ar[r]^-{\partial}\ar[d]^{\iota}
    &
    \mathbf H^{n+1}(\mathfrak U,\mathscr K'^\bullet)
      \ar[d]^{\iota}
    \\
    \mathbf H^n(\mathfrak U',\mathscr K''^\bullet)
      \ar[r]_-{\partial}
    &
    \mathbf H^{n+1}(\mathfrak U',\mathscr K'^\bullet).
  }
\]

Cuando se cumple una de las condiciones precedentes y
$\mathscr K^\bullet\to\mathscr K'^\bullet$ es una equivalencia homotópica (6.1.4), el
isomorfismo correspondiente
$\mathscr H^\bullet(f,\mathscr K^\bullet)\xrightarrow{\sim}
\mathscr H^\bullet(f,\mathscr K'^\bullet)$ es entonces un isomorfismo de
$\partial$-funtores (0, 11.4.4).
\end{env}

\begin{env}[6.2.7]
Todo lo que precede se aplica naturalmente, sin más cambio que el de la
notación, a un complejo $\mathscr K_\bullet$ de $\mathscr O_X$-módulos
cuasicoherentes cuya diferencial es de grado $-1$; basta considerar el
complejo $\mathscr K^\bullet=(\mathscr K^i)$, donde
$\mathscr K^i=\mathscr K_{-i}$ para todo $i\in\mathbf Z$.
\end{env}
